\section{Sistemas con matrices no diagonalizables y exponencial de matrices}

Para entender el comportamiento de los sistemas de ecuaciones diferenciales cuando la matriz de coeficientes no es diagonalizable, analizaremos el caso más sencillo: un bloque de Jordan de orden 2.

Sea el sistema $h' = Ah$, donde la matriz $A \in \mathcal{M}_{2 \times 2}(\mathbb{C})$ tiene la forma:
\[
    A = \begin{pmatrix} \lambda & 0 \\ 1 & \lambda \end{pmatrix}
\]
El espectro de esta matriz es $\sigma(A) = \{\lambda\}$. 

Para encontrar las soluciones del sistema, buscamos un autovector $a$ y un autovector generalizado $b$ que satisfagan las relaciones de la cadena:
\begin{align*}
    (A - \lambda I)a &= 0 \implies Aa = \lambda a \\
    (A - \lambda I)b &= a \implies Ab = \lambda b + a
\end{align*}

Sustituyendo nuestra matriz $A$, la primera ecuación es:
\[
    \begin{pmatrix} 0 & 0 \\ 1 & 0 \end{pmatrix} \begin{pmatrix} a_1 \\ a_2 \end{pmatrix} = \begin{pmatrix} 0 \\ 0 \end{pmatrix} \implies a_1 = 0
\]
Tomamos como autovector $a = \begin{pmatrix} 0 \\ 1 \end{pmatrix}$. El subespacio propio es por tanto $N[A-\lambda I] = L\left[ \begin{pmatrix} 0 \\ 1 \end{pmatrix} \right]$.

Para el autovector generalizado $b = \begin{pmatrix} x \\ y \end{pmatrix}$, resolvemos la segunda ecuación:
\[
    \begin{pmatrix} 0 & 0 \\ 1 & 0 \end{pmatrix} \begin{pmatrix} x \\ y \end{pmatrix} = \begin{pmatrix} 0 \\ 1 \end{pmatrix} \implies x = 1
\]
Podemos tomar $b = \begin{pmatrix} 1 \\ 0 \end{pmatrix}$.

Con estos vectores, construimos una solución generalizada de la forma $h(t) = e^{\lambda t}(ta + b)$. Sustituyendo nuestros vectores obtenidos:
\[
    h(t) = e^{\lambda t} \left[ t \begin{pmatrix} 0 \\ 1 \end{pmatrix} + \begin{pmatrix} 1 \\ 0 \end{pmatrix} \right] = \begin{pmatrix} e^{\lambda t} \\ t e^{\lambda t} \end{pmatrix}
\]

Esto nos permite construir la Matriz Fundamental de Soluciones (M.F.S.), denotada como $\Phi(t)$:
\[
    \Phi(t) = \begin{pmatrix} e^{\lambda t} & 0 \\ t e^{\lambda t} & e^{\lambda t} \end{pmatrix}
\]
Cuyo determinante es $|\Phi(t)| = e^{2\lambda t} \neq 0$, garantizando que las columnas forman un sistema fundamental para la solución general $\mathcal{L}(t)$.

\subsection{Cálculo mediante la exponencial de una matriz}

De manera alternativa, sabemos que la solución al sistema $h' = Ah$ viene dada por la matriz exponencial $e^{At}$, definida mediante su serie de Taylor:
\[
    e^{At} = \sum_{n \ge 0} \frac{(At)^n}{n!}
\]
Para calcularla de forma exacta, descomponemos nuestra matriz $A$ como la suma de su parte diagonal y su parte nilpotente:
\[
    A = \lambda I_2 + N = \begin{pmatrix} \lambda & 0 \\ 0 & \lambda \end{pmatrix} + \begin{pmatrix} 0 & 0 \\ 1 & 0 \end{pmatrix}
\]
Dado que la matriz identidad conmuta con cualquier matriz ($(\lambda I_2)N = N(\lambda I_2)$), podemos aplicar la propiedad de la exponencial para sumas que conmutan ($AB=BA \implies e^{A+B} = e^A e^B$):
\[
    e^{At} = e^{(\lambda I_2 + N)t} = e^{\lambda I_2 t} e^{Nt}
\]

La matriz $e^{\lambda I_2 t}$ es simplemente $e^{\lambda t} I_2$. Para la matriz nilpotente $N$, notamos que $N^2 = 0$, por lo que la serie de Taylor se trunca abruptamente en el primer término:
\[
    e^{Nt} = I_2 + Nt + \frac{(Nt)^2}{2!} + \dots = \begin{pmatrix} 1 & 0 \\ 0 & 1 \end{pmatrix} + \begin{pmatrix} 0 & 0 \\ t & 0 \end{pmatrix} = \begin{pmatrix} 1 & 0 \\ t & 1 \end{pmatrix}
\]
Multiplicando ambos resultados, recuperamos la Matriz Fundamental de Soluciones que obtuvimos por el método anterior:
\[
    e^{At} = \left( e^{\lambda t} I_2 \right) \begin{pmatrix} 1 & 0 \\ t & 1 \end{pmatrix} = \begin{pmatrix} e^{\lambda t} & 0 \\ t e^{\lambda t} & e^{\lambda t} \end{pmatrix}
\]

\subsection{Repaso de Álgebra Lineal: Teorema de Descomposición}

Para extender este proceso a matrices de orden arbitrario, necesitamos recordar los fundamentos teóricos sobre autoespacios generalizados.

Sea $A = (a_{ij})_{1 \le i,j \le N}$ una matriz cuadrada de orden $N$ con coeficientes complejos. Su espectro está formado por las raíces de su polinomio característico:
\[
    \sigma(A) = \{ \lambda : \det(A - \lambda I) = 0 \} = \{\lambda_1, \dots, \lambda_q\} \quad \text{con } \lambda_i \neq \lambda_j \text{ si } i \neq j
\]

Para cada autovalor $\lambda \in \sigma(A)$, definimos su \textbf{autoespacio} (o núcleo) como $N[A-\lambda I]$. Sabemos que la multiplicidad geométrica de un autovalor está acotada por su multiplicidad algebraica:
\[
    \text{Mult. Geométrica} = \dim N[A-\lambda I] \leq m_{alg}(\lambda)
\]

Si evaluamos las potencias sucesivas del operador $(A-\lambda I)$, obtenemos la cadena de inclusiones de los \textbf{autoespacios generalizados}:
\[
    N[A-\lambda I] \subsetneq N[(A-\lambda I)^2] \subsetneq \dots \subsetneq N[(A-\lambda I)^{\nu(\lambda)}]
\]
Esta cadena de inclusiones es estricta hasta llegar a un cierto entero $\nu(\lambda)$, llamado \textbf{índice} o \textbf{ascendente algebraico} del autovalor, a partir del cual el subespacio se estabiliza y ya no crece más.

Esta estructura se consolida en el \textbf{Teorema de Descomposición (Teorema de Jordan)}, el cual establece:
\begin{enumerate}
    \item La multiplicidad algebraica coincide con la dimensión del autoespacio generalizado máximo:
    \[
        m_{alg}(\lambda_j) = \dim N[(A-\lambda_j I)^{\nu(\lambda_j)}]
    \]
    
    \item El espacio vectorial completo $\mathbb{C}^N$ se puede expresar como suma directa de estos autoespacios generalizados:
    \[
        \mathbb{C}^N = \bigoplus_{j=1}^q N[(A-\lambda_j I)^{\nu(\lambda_j)}]
    \]
    
    \item Cada uno de estos subespacios es invariante bajo la transformación de $A$. Es decir, si restringimos la matriz $A$ al subespacio $j$-ésimo (denotándolo como $A_j$), la matriz global $A$ adopta una forma diagonal por bloques conformada por estas submatrices $A_1, A_2, \dots, A_q$.
\end{enumerate}
\subsection{Descomposición y Base de Jordan (Construcción del Cuadro)}

Supongamos que el espectro de la matriz $A$ es $\sigma(A)=\{\lambda_1, \dots, \lambda_s\}$ con autovalores distintos, es decir, $\lambda_i \neq \lambda_j$ si $i \neq j$. 

El espacio vectorial completo se puede expresar como suma directa de los subespacios propios generalizados:
\[
    \mathbb{C}^N = \bigoplus_{j=1}^s N[(A-\lambda_j I)^{\nu(\lambda_j)}]
\]

Nuestro objetivo es construir la base de Jordan en cada uno de estos subespacios $N[(A-\lambda_j I)^{\nu(\lambda_j)}]$. Centrémonos en un autovalor genérico $\lambda$, cuyo índice de estabilización es $\nu$.

Para hallar la base, descomponemos el núcleo máximo estabilizado en una suma directa de las "capas" o diferencias entre los núcleos iterados consecutivos:
\begin{align*}
    N[(A-\lambda I)^\nu] &= \left( N[(A-\lambda I)^\nu] \setminus N[(A-\lambda I)^{\nu-1}] \right) \\
    &\quad \oplus \left( N[(A-\lambda I)^{\nu-1}] \setminus N[(A-\lambda I)^{\nu-2}] \right) \\
    &\quad \oplus \dots \\
    &\quad \oplus \left( N[(A-\lambda I)^2] \setminus N[(A-\lambda I)] \right) \\
    &\quad \oplus N[(A-\lambda I)]
\end{align*}

\subsubsection*{El algoritmo: Construcción del cuadro de cadenas}

Vamos a rellenar cada sumando (cada capa) de arriba hacia abajo. Esto se representa mediante una tabla donde cada fila es un nivel ("sumando") y cada columna es una "cadena" de vectores generada por el $(A-\lambda I)$.

\begin{itemize}
    \item \textbf{1º Sumando (Nivel $\nu$):} Tomamos los vectores que faltan para rellenar la dimensión de la capa más alta. Los llamamos $v_{h_{\nu-1}+1}, \dots, v_{h_\nu}$.
    \item \textbf{2º Sumando (Nivel $\nu-1$):} Primero, bajamos los vectores del nivel anterior aplicando el operador $(A-\lambda I)$. Después, si la dimensión de esta capa lo requiere, añadimos nuevos vectores independientes para rellenar este nivel: $v_{h_{\nu-2}+1}, \dots, v_{h_{\nu-1}}$.
    \item Repetimos el proceso aplicando el ascensor $(A-\lambda I)$ a \textit{todos} los vectores que ya tenemos, bajándolos al siguiente nivel y añadiendo nuevos solo si hace falta.
    \item \textbf{Último Sumando (Nivel 1):} Llegamos a la capa de autovectores puros. Las cadenas terminan aquí.
\end{itemize}

El esquema visual (Cuadro de Jordan) queda así:

\begin{center}
\renewcommand{\arraystretch}{1.8}
\begin{tabular}{l c c c}
\textbf{1º sumando:} & $v_{h_{\nu-1}+1}, \dots, v_{h_\nu}$ & & \\
& $\downarrow \scriptstyle{(A-\lambda I)}$ & & \\
\textbf{2º sumando:} & $(A-\lambda I)v_{h_{\nu-1}+1}, \dots$ & $v_{h_{\nu-2}+1}, \dots, v_{h_{\nu-1}}$ & \\
& $\downarrow \scriptstyle{(A-\lambda I)}$ & $\downarrow \scriptstyle{(A-\lambda I)}$ & \\
\textbf{3º sumando:} & $(A-\lambda I)^2 v_{h_{\nu-1}+1}, \dots$ & $(A-\lambda I)v_{h_{\nu-2}+1}, \dots$ & $v_{h_{\nu-3}+1}, \dots$ \\
$\vdots$ & $\vdots$ & $\vdots$ & $\vdots$ \\
\textbf{Último sumando:} & $(A-\lambda I)^{\nu-1}v_{h_{\nu-1}+1}, \dots$ & $(A-\lambda I)^{\nu-2}v_{h_{\nu-2}+1}, \dots$ & $\dots \quad v_1, \dots, v_{h_1}$
\end{tabular}
\end{center}

Al final de este proceso, \textbf{hemos completado la base de Jordan}. Cada columna de esta tabla es una "cadena" que generará un bloque (o caja) en la matriz de Jordan. La longitud de la columna nos dirá el tamaño de esa caja.

Tomemos una de estas cadenas generada a partir de un vector cabeza:
\begin{align*}
    w_1 &= v_1 \\
    w_2 &= (A-\lambda I)w_1 \\
    w_3 &= (A-\lambda I)w_2 = (A-\lambda I)^2 w_1 \\
    w_4 &= (A-\lambda I)w_3 \\
    &\vdots \\
    w_\nu &= (A-\lambda I)w_{\nu-1}
\end{align*}

Despejando la acción de la matriz $A$ sobre los vectores de esta cadena, obtenemos:
\begin{align*}
    A w_1 &= \lambda w_1 + w_2 \\
    A w_2 &= \lambda w_2 + w_3 \\
    A w_3 &= \lambda w_3 + w_4 \\
    &\vdots \\
    A w_{\nu-1} &= \lambda w_{\nu-1} + w_\nu \\
    A w_\nu &= \lambda w_\nu \quad \text{(autovector)}
\end{align*}

Matricialmente, la actuación de $A$ sobre esta base $\{w_1, w_2, \dots, w_\nu\}$ da lugar a una caja de Jordan $J_{\lambda, \nu}$:
\[
    J_{\lambda, \nu} = \begin{pmatrix} 
    \lambda & 0 & 0 & \dots & 0 \\ 
    1 & \lambda & 0 & \dots & 0 \\ 
    0 & 1 & \lambda & \dots & 0 \\ 
    \vdots & \vdots & \ddots & \ddots & \vdots \\ 
    0 & 0 & \dots & 1 & \lambda 
    \end{pmatrix}
\]

El número de cajas de Jordan asociadas a cada sumando (nivel) se deduce de las diferencias de dimensiones de los núcleos:
\begin{itemize}
    \item \textbf{1º sumando:} $h_\nu - h_{\nu-1}$ cajas de Jordan $J_{\lambda, \nu}$
    \item \textbf{2º sumando:} $h_{\nu-1} - h_{\nu-2}$ cajas de Jordan $J_{\lambda, \nu-1}$
    \item \textbf{3º sumando:} $h_{\nu-2} - h_{\nu-3}$ cajas de Jordan $J_{\lambda, \nu-2}$
    \item \dots
    \item \textbf{Penúltimo sumando:} $h_2 - h_1$ cajas de Jordan $J_{\lambda, 2}$
    \item \textbf{Último sumando:} $h_1 - h_0$ cajas de Jordan $J_{\lambda, 1}$
\end{itemize}

Las cajas de Jordan tienen la siguiente estructura (con los unos en la subdiagonal):
\[
    J_{\lambda, 1} = (\lambda), \quad 
    J_{\lambda, 2} = \begin{pmatrix} \lambda & 0 \\ 1 & \lambda \end{pmatrix}, \quad 
    J_{\lambda, 3} = \begin{pmatrix} \lambda & 0 & 0 \\ 1 & \lambda & 0 \\ 0 & 1 & \lambda \end{pmatrix}, \quad 
    J_{\lambda, 4} = \begin{pmatrix} \lambda & 0 & 0 & 0 \\ 1 & \lambda & 0 & 0 \\ 0 & 1 & \lambda & 0 \\ 0 & 0 & 1 & \lambda \end{pmatrix}
\]

Para cada autovalor $\lambda$, la matriz de Jordan asociada se construye colocando estos bloques en la diagonal principal. El número de bloques de cada tamaño viene dado por las diferencias de dimensiones de los núcleos iterados ($h_k$). Visualmente, la matriz $J$ asociada a $\lambda$ queda así:

\[
J = \begin{pmatrix}
    \begin{matrix} J_{\lambda,\nu} \\ & \ddots \\ & & J_{\lambda,\nu} \end{matrix} & & & & \\
    & \begin{matrix} J_{\lambda,\nu-1} \\ & \ddots \\ & & J_{\lambda,\nu-1} \end{matrix} & & & \\
    & & \ddots & & \\
    & & & \begin{matrix} J_{\lambda,2} \\ & \ddots \\ & & J_{\lambda,2} \end{matrix} & \\
    & & & & \begin{matrix} J_{\lambda,1} \\ & \ddots \\ & & J_{\lambda,1} \end{matrix}
\end{pmatrix}
\begin{matrix}
    \left.\begin{matrix} \vphantom{J_{\lambda,\nu}} \\ \vphantom{\ddots} \\ \vphantom{J_{\lambda,\nu}} \end{matrix}\right\} & h_\nu - h_{\nu-1} \text{ cajas} \\
    \left.\begin{matrix} \vphantom{J_{\lambda,\nu-1}} \\ \vphantom{\ddots} \\ \vphantom{J_{\lambda,\nu-1}} \end{matrix}\right\} & h_{\nu-1} - h_{\nu-2} \text{ cajas} \\
    \vphantom{\ddots} & \vdots \\
    \left.\begin{matrix} \vphantom{J_{\lambda,2}} \\ \vphantom{\ddots} \\ \vphantom{J_{\lambda,2}} \end{matrix}\right\} & h_2 - h_1 \text{ cajas} \\
    \left.\begin{matrix} \vphantom{J_{\lambda,1}} \\ \vphantom{\ddots} \\ \vphantom{J_{\lambda,1}} \end{matrix}\right\} & h_1 - h_0 \text{ cajas}
\end{matrix}
\]

\vspace{0.5cm}

\ejemplo{
    Sea $A$ una matriz con espectro $\sigma(A) = \{\alpha, \beta, \gamma\}$, donde $\alpha \neq \beta$, $\beta \neq \gamma$ y $\alpha \neq \gamma$. Se conocen las siguientes dimensiones de sus núcleos iterados:
    
    \begin{itemize}
        \item Para el autovalor $\alpha$:
        \[ \dim N[(A-\alpha I)] = 3, \quad \dim N[(A-\alpha I)^2] = 3 \]
        
        \item Para el autovalor $\beta$:
        \[ \dim N[(A-\beta I)] = 3, \quad \dim N[(A-\beta I)^2] = 4 \]
        \[ \dim N[(A-\beta I)^3] = 5, \quad \dim N[(A-\beta I)^4] = 5 \]
        
        \item Para el autovalor $\gamma$:
        \[ \dim N[(A-\gamma I)] = 2, \quad \dim N[(A-\gamma I)^2] = 3 \]
        \[ \dim N[(A-\gamma I)^3] = 3 \]
    \end{itemize}
    
    Se pide:
    \begin{enumerate}
        \item Orden de A.
        \item Multiplicidad algebraica de $\alpha, \beta, \gamma$.
        \item Forma canónica de Jordan de A.
    \end{enumerate}
    \textbf{Índices:}
    \begin{align*}
        \nu(\alpha) = 1 &\implies m_{alg}(\alpha) = \dim N[(A-\alpha I)^1] = 3 \\
        \nu(\beta) = 3 &\implies m_{alg}(\beta) = \dim N[(A-\beta I)^3] = 5 \\
        \nu(\gamma) = 2 &\implies m_{alg}(\gamma) = \dim N[(A-\gamma I)^2] = 3
    \end{align*}

    \textbf{Suma directa:}
    \[
        \mathbb{C}^{11} = \underbrace{N[(A-\alpha I)]}_{\dim 3} \oplus \underbrace{N[(A-\beta I)^3]}_{\dim 5} \oplus \underbrace{N[(A-\gamma I)^2]}_{\dim 3}
    \]

    \textbf{Orden de A:}
    \[
        N = 3 + 5 + 3 = 11
    \]

    Vamos a calcular la base de Jordan:

    \textbf{1) Para el autovalor $\alpha$:}
    Analizamos el núcleo $N[(A-\alpha I)]$:
    \[
        \begin{matrix}
            v_{\alpha,1} \\
            v_{\alpha,2} \\
            v_{\alpha,3} 
        \end{matrix}
        \quad (A-\alpha I)v_{\alpha,i} = 0 \implies A v_{\alpha,i} = \alpha v_{\alpha,i}
    \]
    Estos tres vectores son autovectores y generarán tres cajas $J_{\alpha,1}$.

    \textbf{2) Para el autovalor $\beta$:}
    Analizamos el núcleo de índice máximo:
    \[
        N[(A-\beta I)^3] = \left( N[(A-\beta I)^3] \setminus N[(A-\beta I)^2] \right) \oplus \left( N[(A-\beta I)^2] \setminus N[(A-\beta I)] \right) \oplus N[(A-\beta I)]
    \]

    Construimos la cadena de forma descendente:
    \begin{align*}
        v_{\beta,2} &= (A-\beta I)v_{\beta,1} \\
        v_{\beta,3} &= (A-\beta I)v_{\beta,2} = (A-\beta I)^2 v_{\beta,1}
    \end{align*}

    Despejando la acción de $A$ sobre la cadena:
    \begin{align*}
        A v_{\beta,1} &= \beta v_{\beta,1} + v_{\beta,2} \\
        A v_{\beta,2} &= \beta v_{\beta,2} + v_{\beta,3} \\
        A v_{\beta,3} &= \beta v_{\beta,3}
    \end{align*}

    Esto da lugar a la primera caja de orden 3:
    \[
        \implies \begin{pmatrix} \beta & 0 & 0 \\ 1 & \beta & 0 \\ 0 & 1 & \beta \end{pmatrix} = J_{\beta,3}
    \]

    Ahora completamos bases. No hay que añadir nada en el $2º$ sumando. En el autoespacio completamos con $v_{\beta,4}$ y $v_{\beta,5}$, lo que aportará dos cajas $J_{\beta,1}$.

    \textbf{Matriz de Jordan $J$ final:}
    Agrupando los bloques para $\alpha$, $\beta$ (y asumiendo el desarrollo análogo para $\gamma$), la matriz completa en la base de Jordan es:

    \[
    J = \left(
    \begin{array}{ccc|ccccc|ccc}
    \alpha & 0 & 0 & 0 & 0 & 0 & 0 & 0 & 0 & 0 & 0 \\
    0 & \alpha & 0 & 0 & 0 & 0 & 0 & 0 & 0 & 0 & 0 \\
    0 & 0 & \alpha & 0 & 0 & 0 & 0 & 0 & 0 & 0 & 0 \\
    \hline
    0 & 0 & 0 & \beta & 0 & 0 & 0 & 0 & 0 & 0 & 0 \\
    0 & 0 & 0 & 1 & \beta & 0 & 0 & 0 & 0 & 0 & 0 \\
    0 & 0 & 0 & 0 & 1 & \beta & 0 & 0 & 0 & 0 & 0 \\
    0 & 0 & 0 & 0 & 0 & 0 & \beta & 0 & 0 & 0 & 0 \\
    0 & 0 & 0 & 0 & 0 & 0 & 0 & \beta & 0 & 0 & 0 \\
    \hline
    0 & 0 & 0 & 0 & 0 & 0 & 0 & 0 & \gamma & 0 & 0 \\
    0 & 0 & 0 & 0 & 0 & 0 & 0 & 0 & 1 & \gamma & 0 \\
    0 & 0 & 0 & 0 & 0 & 0 & 0 & 0 & 0 & 0 & \gamma
    \end{array}
    \right)
    \]
}


\subsection{Forma Canónica Real}

\subsubsection{Introducción y Propiedades Fundamentales}

Consideremos una matriz cuadrada con coeficientes reales $A = (a_{ij})_{1 \le i,j \le N} \in \mathcal{M}_{N \times N}(\mathbb{R})$. 

Su polinomio característico, denotado como $P(z) = \det(A - zI)$, es un polinomio con coeficientes estrictamente reales. Por el Teorema Fundamental del Álgebra y las propiedades de la conjugación compleja, sabemos que si un número complejo $\lambda$ es raíz del polinomio, su conjugado $\overline{\lambda}$ también lo es, y ambos comparten la misma multiplicidad algebraica:
$$ P(\lambda) = 0 \iff P(\overline{\lambda}) = 0, \quad \text{siendo } m_{alg}(\lambda) = m_{alg}(\overline{\lambda}) $$

Esta simetría no se limita a los autovalores, sino que se traslada directamente a los espacios de autovectores y autovectores generalizados. Para cualquier entero $k \ge 1$, las ecuaciones que definen el núcleo se conjugan de forma natural, ya que la matriz $A$ es real ($\overline{A} = A$):
$$ (A - \lambda I)^k u = 0 \iff (A - \overline{\lambda} I)^k \overline{u} = 0 $$

Por lo tanto, un vector pertenece al núcleo generalizado de $\lambda$ si y solo si su conjugado pertenece al núcleo generalizado de $\overline{\lambda}$:
$$ u \in \ker(A - \lambda I)^k \iff \overline{u} \in \ker(A - \overline{\lambda} I)^k $$
Lo cual implica que los subespacios son mutuamente conjugados:
$$ \ker(A - \overline{\lambda} I)^k = \overline{\ker(A - \lambda I)^k} $$

\subsubsection{Descomposición del Espacio y Bases Reales}

Por el Teorema de Descomposición Primaria, el espacio vectorial complejo $\mathbb{C}^N$ se puede expresar como suma directa de los núcleos generalizados asociados a cada autovalor. Si separamos explícitamente un autovalor complejo $\lambda$ y su conjugado $\overline{\lambda}$ del resto del espectro, tenemos:
$$ \mathbb{C}^N = \ker(A - \lambda_1 I)^{\nu_1} \oplus \dots \oplus \ker(A - \lambda I)^{\nu(\lambda)} \oplus \ker(A - \overline{\lambda} I)^{\nu(\overline{\lambda})} $$

Donde $\nu(\lambda)$ representa el índice del autovalor (el tamaño del bloque de Jordan más grande asociado a $\lambda$). Si hallamos una base de autovectores generalizados para el subespacio asociado a $\lambda$:
$$ \mathcal{B}_\lambda = \{u_{\lambda,1}, u_{\lambda,2}, \dots, u_{\lambda,m_{\lambda}}\} $$
Automáticamente, conjugando cada elemento, obtenemos una base para el subespacio de $\overline{\lambda}$:
$$ \mathcal{B}_{\overline{\lambda}} = \{\overline{u_{\lambda,1}}, \overline{u_{\lambda,2}}, \dots, \overline{u_{\lambda,m_{\lambda}}}\} $$

\textbf{Transición a la base real:} La unión de $\mathcal{B}_\lambda$ y $\mathcal{B}_{\overline{\lambda}}$ proporciona una base compleja para la suma directa de ambos subespacios. Sin embargo, para obtener una matriz semejante con coeficientes puramente reales, debemos transformar esta base compleja en una \textbf{base real}. Esto se logra extrayendo las partes reales e imaginarias de los vectores de las cadenas de Jordan, usando las combinaciones:
$$ u + \overline{u} = 2\Re(u), \quad i(u - \overline{u}) = -2\Im(u) $$

Definimos la base real equivalente agrupando las partes reales e imaginarias. Para cada vector $u_{\lambda, j}$, introducimos el par $\{\Re(u_{\lambda, j}), \Im(u_{\lambda, j})\}$, de modo que el subespacio generado por el par complejo $\{u, \overline{u}\}$ es exactamente el mismo que el generado por el par real $\{\Re(u), \Im(u)\}$.

\subsubsection{Estructura del Bloque de Jordan Real}

Antes de ilustrar el ejemplo, es conveniente definir el bloque estándar de rotación-dilatación. Si $\lambda = \alpha + i\beta$ con $\beta \neq 0$, la acción de $A$ sobre un autovector complejo $v = \Re(v) + i\Im(v)$ se traduce en un bloque real $2 \times 2$ de la forma:
$$ D_{\lambda} = \begin{pmatrix} \alpha & -\beta \\ \beta & \alpha \end{pmatrix} $$
Un \textbf{Bloque de Jordan Real} de tamaño $2k \times 2k$ asociado a $\lambda = \alpha + i\beta$ tiene la forma de una matriz bidiagonal por bloques:
$$ J_{\mathbb{R}}(\lambda, k) = \begin{pmatrix}
D_{\lambda} & 0 & \cdots & 0 \\
I_2 & D_{\lambda} & \cdots & 0 \\
\vdots & \ddots & \ddots & \vdots \\
0 & \cdots & I_2 & D_{\lambda}
\end{pmatrix} $$
Donde $I_2$ es la matriz identidad de tamaño $2 \times 2$.

\ejemplo{
    Supongamos que una matriz real $A$ tiene un espectro puramente complejo formado por un par conjugado:
    $$ \sigma(A) = \{\lambda, \overline{\lambda}\}, \quad \text{donde } \lambda = \alpha + i\beta, \; (\beta > 0) \quad \text{y} \quad \overline{\lambda} = \alpha - i\beta $$

    Analizamos las dimensiones de sus núcleos para determinar la estructura. Supongamos las siguientes dimensiones:
    $$ \dim(\ker(A - \lambda I)) = 1 $$
    $$ \dim(\ker(A - \lambda I)^2) = \dim(\ker(A - \lambda I)^3) = 2 $$

    Esto indica que las cadenas se estabilizan en el paso 2, por lo que el índice es $\nu(\lambda) = 2$. La multiplicidad algebraica es $m_{alg}(\lambda) = m_{alg}(\overline{\lambda}) = 2$.

    \paragraph{Construcción de la Base Compleja.}
    Buscamos un vector que esté en el núcleo de orden 2 pero no en el de orden 1:
    $$ v_{\lambda,1} \in \ker((A - \lambda I)^2) \setminus \ker(A - \lambda I) $$
    Generamos la cadena de Jordan aplicando la transformación:
    $$ (A - \lambda I) v_{\lambda,1} = v_{\lambda,2} $$
    Como $v_{\lambda,2} \in \ker(A - \lambda I)$, es un autovector propio ($Av_{\lambda,2} = \lambda v_{\lambda,2}$). La base para $\lambda$ es $\mathcal{B} = \{v_{\lambda,1}, v_{\lambda,2}\}$. 

    Por conjugación, la base para $\overline{\lambda}$ es $\overline{\mathcal{B}} = \{v_{\overline{\lambda},1}, v_{\overline{\lambda},2}\}$, donde $v_{\overline{\lambda},1} = \overline{v_{\lambda,1}}$ y $v_{\overline{\lambda},2} = \overline{v_{\lambda,2}}$.

    En la base compleja total $L = \{v_{\lambda,1}, v_{\lambda,2}, v_{\overline{\lambda},1}, v_{\overline{\lambda},2}\}$, la matriz adopta su forma canónica de Jordan compleja:
    $$ A |_{L} = \begin{pmatrix}
        J_{\lambda,2} & 0 \\
        0 & J_{\overline{\lambda},2}
    \end{pmatrix} = \begin{pmatrix}
        \lambda & 0 & 0 & 0 \\
        1 & \lambda & 0 & 0 \\
        0 & 0 & \overline{\lambda} & 0 \\
        0 & 0 & 1 & \overline{\lambda}
    \end{pmatrix} $$

    \paragraph{Paso a la Forma Canónica Real.}
    A partir de la cadena de Jordan, tenemos el sistema:
    $$ \begin{cases}
        A v_{\lambda,1} = \lambda v_{\lambda,1} + v_{\lambda,2} \\
        A v_{\lambda,2} = \lambda v_{\lambda,2}
    \end{cases} $$

    Sustituimos $v_{\lambda, j} = \Re(v_{\lambda, j}) + i \Im(v_{\lambda, j})$ y separamos la parte real de la imaginaria. 
    Para el autovector ($v_2$):
    $$ A(\Re(v_2) + i \Im(v_2)) = (\alpha + i\beta)(\Re(v_2) + i \Im(v_2)) = (\alpha \Re(v_2) - \beta \Im(v_2)) + i(\beta \Re(v_2) + \alpha \Im(v_2)) $$
    Igualando componentes:
    $$ A \Re(v_2) = \alpha \Re(v_2) - \beta \Im(v_2) $$
    $$ A (-\Im(v_2)) = \beta \Re(v_2) + \alpha (-\Im(v_2)) $$

    Repitiendo el proceso para el vector generalizado ($v_1$):
    $$ A \Re(v_1) = \alpha \Re(v_1) - \beta \Im(v_1) + \Re(v_2) $$
    $$ A (-\Im(v_1)) = \beta \Re(v_1) + \alpha (-\Im(v_1)) + (-\Im(v_2)) $$

    Definimos nuestra base real ordenada:
    $$ \mathcal{B}_{\mathbb{R}} = \{\Re(v_{\lambda,1}), -\Im(v_{\lambda,1}), \Re(v_{\lambda,2}), -\Im(v_{\lambda,2})\} $$

    Si escribimos los coeficientes de estas ecuaciones por columnas, la matriz $A$ restringida a esta base adopta su Forma Canónica Real definitiva:
    $$ A |_{\mathcal{B}_{\mathbb{R}}} = \begin{pmatrix}
        \alpha & -\beta & 0 & 0 \\
        \beta & \alpha & 0 & 0 \\
        1 & 0 & \alpha & -\beta \\
        0 & 1 & \beta & \alpha
    \end{pmatrix} $$
    Esta matriz revela de forma explícita la estructura de dilatación ($\alpha$), rotación ($\beta$) y acoplamiento nilpotente ($I_2$) del sistema original.
}


\begin{teorema}
    Para cada matriz $A \in \mathcal{M}_{N \times N}(\mathbb{C})$, la sucesión de sumas parciales 
    \[
        S_n (A) = \sum_{k=0}^n \frac{A^k}{k!}, \quad n \ge 0
    \]
    es convergente en $\mathcal{M}_{N \times N}(\mathbb{C})$ y su límite se denomina \textbf{exponencial de la matriz $A$}, denotado como $e^A$:
    \[
        e^A = \lim_{n \to \infty} S_n(A) = \sum_{k=0}^\infty \frac{A^k}{k!}
    \]
\end{teorema}

\begin{proof}
    Tomamos la norma operador
    \[
        \|T\|_{\mathcal{L}(\mathbb{C}^N)} = \max_{\|x\|=1} \|Tx\|
    \]
    que cumple $\|T^2x\| \le \|T\| \|Tx\|$ y, por inducción, $\|T^k x\| \le \|T\|^k \|x\|$ para todo $k \ge 1$.
    Tenemos que $\|x\| = 1$, por lo que $\|T^k x\| \le \|T\|^k$. Por lo tanto,
    \[
        \|T^k\| = \max_{\|x\|=1} \|T^k x\| \le \|T\|^k
    \]
    Queremos mostrar que $\{S_n(A)\}$ es una sucesión de Cauchy. 
    \[
        \|S_{n+m} (A) - S_n(A)\|_{\mathcal{L}(\mathbb{C}^N)} = \left\| \sum_{k=n+1}^{n+m} \frac{A^k}{k!} \right\|_{\mathcal{L}(\mathbb{C}^N)} \le \sum_{k=n+1}^{n+m} \frac{\|A^k\|_{\mathcal{L}(\mathbb{C}^N)}}{k!} \le \sum_{k=n+1}^{n+m} \frac{\|A\|_{\mathcal{L}(\mathbb{C}^N)}^k}{k!}
    \]
    Pero sabemos que la siguiente serie es convergente:
    \[
        e^{\|A\|_{\mathcal{L}(\mathbb{C}^N)}} = \sum_{k=0}^\infty \frac{\|A\|_{\mathcal{L}(\mathbb{C}^N)}^k}{k!} < \infty
    \]
    luego $\forall \varepsilon > 0$, $\exists N_\varepsilon$ tal que $\sum_{k=n+1}^\infty \frac{\|A\|_{\mathcal{L}(\mathbb{C}^N)}^k}{k!} < \varepsilon$ para todo $n \ge N_\varepsilon$. Por lo tanto, $\{S_n(A)\}$ es una sucesión de Cauchy y converge a un límite que denominamos $e^A$.
\end{proof}


\paragraph{Propiedades:}

\begin{itemize}
    \item Si $B = P^{-1} A P$ con $P$ invertible, entonces $e^B = P^{-1} e^A P$. Es decir, la función exponencial es invariante por semejanza. \\
    Aplicación: $J = P^{-1} A P \quad A = P J P^{-1} \implies e^A = P e^J P^{-1}$, lo que nos permite calcular $e^A$ a partir de la forma de Jordan de $A$.
    \[
        S_n(B) = \sum_{k=0}^n \frac{B^k}{k!} = \sum_{k=0}^n \frac{(P^{-1} A P)^k}{k!} = \sum_{k=0}^n P^{-1} \frac{A^k}{k!} P = P^{-1} \sum_{k=0}^n \frac{A^k}{k!} P = P^{-1} S_n(A) P
    \]
    \item $\|e^A\| \leq e^{\|A\|}$, lo que garantiza que la función exponencial de una matriz es siempre una matriz acotada.
    \[
        S_n(A) = \sum_{k=0}^n \frac{A^k}{k!} \implies \|S_n(A)\| = \left\| \sum_{k=0}^n \frac{A^k}{k!} \right\| \leq \sum_{k=0}^n \frac{\|A^k\|}{k!} \leq \sum_{k=0}^n \frac{\|A\|^k}{k!} \leq \sum_{k=0}^\infty \frac{\|A\|^k}{k!} = e^{\|A\|}
    \]
    \item Si dos matrices $A$ y $B$ conmutan ($AB = BA$), entonces $e^{A+B} = e^A e^B$. \\
    \[
        (A+B)^2 = A^2 + AB + BA + B^2 = A^2 + 2AB + B^2, \quad \text{si} \quad AB = BA
    \]
    \[
        (A+B)^3 = (A+B)(A+B)^2 = (A+B)(A^2 + 2AB + B^2) = A^3 + 3A^2B + 3AB^2 + B^3, \quad \text{si} \quad AB = BA
    \]
    \[
        (A+B)^n = \sum_{k=0}^n \binom{n}{k} A^k B^{n-k}, \quad \text{si} \quad AB = BA
    \]
    Por lo tanto,
    \[
        e^{A+B} = \sum_{n=0}^\infty \frac{(A+B)^n}{n!} = \sum_{n=0}^\infty \frac{1}{n!} \sum_{k=0}^n \binom{n}{k} A^k B^{n-k} = \sum_{n=0}^\infty \sum_{k=0}^n \frac{A^k}{k!} \frac{B^{n-k}}{(n-k)!} = \left( \sum_{k=0}^\infty \frac{A^k}{k!} \right) \left( \sum_{m=0}^\infty \frac{B^m}{m!} \right) = e^A e^B
    \]
    \item $e^0 = I$, donde $0$ es la matriz nula y $I$ es la matriz identidad. \\
    \[
        e^0 = \sum_{k=0}^\infty \frac{0^k}{k!} = \frac{0^0}{0!} + \sum_{k=1}^\infty \frac{0^k}{k!} = I + 0 = I
    \]
    \item  
    \[
        \frac{d}{dt} e^{tA} = \frac{d}{dt} \sum_{k=0}^\infty \frac{(tA)^k}{k!} = \sum_{k=1}^\infty \frac{k(tA)^{k-1} A}{k!} = A \sum_{k=1}^\infty \frac{(tA)^{k-1}}{(k-1)!} = A e^{tA}
    \]
    Luego $e^{tA}$ es solución del sistema $u' = A u$ con condición inicial $u(0) = I$. Luego es una matriz fundamental de soluciones porque tiene determinante distinto de cero para todo $t$. En particular, 
    \[
        (e^{tA})^{-1} = e^{-tA}
    \]
\end{itemize}

\[
    e^{\begin{pmatrix}
        \lambda & 0 \\
        1 & \lambda
    \end{pmatrix}} = e^{\lambda I_2 + \begin{pmatrix}
        0 & 0 \\
        1 & 0
    \end{pmatrix}} = e^{\lambda I_2} e^{\begin{pmatrix}
        0 & 0 \\
        1 & 0
    \end{pmatrix}} = e^{\lambda I_2} \sum_{k=0}^\infty \frac{1}{k!} \begin{pmatrix}
        1 & 0 \\
        1 & 1
    \end{pmatrix}^k = \sum_{k=0}^\infty \frac{(\lambda I_2)^k}{k!} \begin{pmatrix}
        1 & 0 \\
        1 & 1
    \end{pmatrix}^k = \begin{pmatrix}
        e^{\lambda} & 0 \\
        e^{\lambda} & e^{\lambda}
    \end{pmatrix}
\]

\begin{teorema}
    Para cada matriz $A \in \mathcal{M}_{N \times N}(\mathbb{C})$, la aplicación $\varphi(z) : \mathbb{C} \to \mathcal{M}_{N \times N}(\mathbb{C})$ definida por $\varphi(z) = e^{zA}$ es una función entera, es decir, es holomorfa en todo el plano complejo.\\
    Es decir, $\forall z \in \mathbb{C}$, existe
    \[
        \varphi'(z) = \lim_{h \to 0} \frac{\varphi(z+h) - \varphi(z)}{h}
    \]
    además, $\varphi'(z) = A e^{zA}$ para todo $z \in \mathbb{C}$.
\end{teorema}

\begin{proof}
    Tomamos el cociente incremental:
    \[
        \frac{\varphi(z+h) - \varphi(z)}{h} = \frac{e^{(z+h)A} - e^{zA}}{h} = \frac{e^{zA} e^{hA} - e^{zA}}{h} = e^{zA} \frac{e^{hA} - I}{h} \xrightarrow{h \to 0} e^{zA} A
    \]
    Hemos de probar que $\lim_{h \to 0} \frac{e^{hA} - I}{h} = A$. Para ello, consideramos la serie de potencias de $e^{hA}$:
    \[
        e^{hA} = \sum_{k=0}^\infty \frac{(hA)^k}{k!} = I + hA + \sum_{k=2}^\infty \frac{(hA)^k}{k!}
    \]
    \[
        \| \frac{e^{hA} - I}{h} - A \| = \left\| h \sum_{k=2}^{\infty} h^{k-2} \frac{A^k}{k!} \right\| \le |h| \sum_{k=2}^{\infty} |h|^{k-2} \frac{\|A\|^k}{k!}
    \]

    Si $|h| \le 1$, entonces $|h|^{k-2} \le 1$ para todo $k \ge 2$, de modo que
    \[
        \| \frac{e^{hA} - I}{h} - A \| \le |h| \sum_{k=2}^{\infty} \frac{\|A\|^k}{k!} = |h| (e^{\|A\|} - 1 - \|A\|) \xrightarrow{h \to 0} 0.
    \]
    Por tanto,
    \[
        \lim_{h \to 0} \frac{e^{hA} - I}{h} = A,
    \]
    y así
    \[
        \varphi'(z) = e^{zA}A.
    \]
    Como $e^{zA}$ es serie de potencias en $A$, conmuta con $A$; por ello $e^{zA}A = Ae^{zA}$.

\end{proof}

\begin{teorema}
    Para cada $u_0 \in \mathbb{C}^N$ y $t_0 \in \mathbb{R}$, la única solución del problema de valor inicial
    \[
        \begin{cases}
            u' = Au + B(t) \\
            u(t_0) = u_0
        \end{cases}
    \]
    donde $B : \mathbb{R} \to \mathbb{C}^N$ es continua, ($[\alpha,\beta] \rightarrow \mathbb{C}, t_0 \in [\alpha,\beta]$) es:
    \[
        u(t) = e^{(t-t_0)A} u_0 + \int_{t_0}^t e^{(t-s)A} B(s) \, ds
    \]
\end{teorema}

\begin{proof}
    Sabemos que $e^{tA} = \Phi(t)$ es la Matriz Fundamental de Soluciones (M.F.S.) de $h' = Ah$.
    
    Proponemos $u(t) = e^{tA} v(t)$ (Variación de coeficientes).
    Sustituyendo en la ecuación $u' = Au + B(t)$:
    \[
        A \cdot e^{tA} v(t) + e^{tA} \cdot v'(t) = A \cdot e^{tA} \cdot v(t) + B(t)
    \]
    Cancelando términos nos queda:
    \[
        e^{tA} \cdot v'(t) = B(t) \implies v'(t) = (e^{tA})^{-1} \cdot B(t) = e^{-tA} \cdot B(t)
    \]
    Y la condición inicial nos da: $u(t_0) = e^{t_0 A} v(t_0) = u_0 \implies v(t_0) = e^{-t_0 A} u_0$.
    
    Integrando:
    \[
        v(t) = v(t_0) + \int_{t_0}^t v'(s) \, ds = e^{-t_0 A} u_0 + \int_{t_0}^t e^{-sA} B(s) \, ds
    \]
    Por tanto, multiplicando por $e^{tA}$:
    \[
        u(t) = e^{tA} v(t) = e^{tA} \left( e^{-t_0 A} u_0 + \int_{t_0}^t e^{-sA} B(s) \, ds \right) = e^{(t-t_0)A} u_0 + \int_{t_0}^t e^{(t-s)A} B(s) \, ds
    \]
\end{proof}

\subsection*{Cálculo de exponenciales}

Usando la forma de Jordan $A = P J P^{-1}$ (donde las columnas de $P$ son los autovectores/autovectores generalizados):
\[
    tA = P \cdot (tJ) \cdot P^{-1} \implies e^{tA} = P \cdot e^{tJ} \cdot P^{-1}
\]

Sabemos que $J$ es diagonal por bloques:
\[
    J = \begin{pmatrix} J_1 & & \\ & \ddots & \\ & & J_k \end{pmatrix}, \quad J_i \in \{J_{\lambda,n}, \lambda \in \sigma(A), n \in \{1, \dots, \nu(\lambda)\}\}
\]
Por definición de la exponencial por serie de Taylor:
\[
    e^{tJ} = \sum_{k=0}^{\infty} \frac{(tJ)^k}{k!} = \sum_{k=0}^{\infty} \frac{t^k}{k!} \begin{pmatrix} J_1^k & & \\ & \ddots & \\ & & J_k^k \end{pmatrix} = \begin{pmatrix} e^{tJ_1} & & \\ & \ddots & \\ & & e^{tJ_k} \end{pmatrix}
\]

Para una caja genérica $\overline{J}_{\lambda,n}$ de orden $n$:
\[
    \overline{J}_{\lambda,n} = \begin{pmatrix} 
    \lambda & 0 & \dots & 0 \\ 
    1 & \lambda & \dots & 0 \\
    \vdots & \ddots & \ddots & \vdots \\
    0 & \dots & 1 & \lambda 
    \end{pmatrix} = \lambda \cdot I_n + N_n
\]
Donde $N_n$ es la matriz nilpotente (unos en la subdiagonal). Notemos que sus potencias van desplazando la diagonal de unos hacia abajo, hasta que $N_n^n = 0$ (nilpotente).

Como $(\lambda I_n) N_n = N_n (\lambda I_n)$, conmutan:
\[
    t \cdot J_{\lambda,n} = t \cdot \lambda \cdot I_n + t \cdot N_n \implies e^{t J_{\lambda,n}} = e^{t \lambda I_n + t N_n} = e^{t \lambda I_n} \cdot e^{t N_n}
\]
Desarrollando la serie (que se trunca al llegar a la potencia $n$):
\[
    e^{t J_{\lambda,n}} = e^{\lambda t} \sum_{k=0}^{\infty} \frac{(t N_n)^k}{k!} = e^{\lambda t} \left[ I_n + t N_n + \frac{t^2}{2!} N_n^2 + \frac{t^3}{3!} N_n^3 + \dots + \frac{t^{n-1}}{(n-1)!} N_n^{n-1} \right]
\]

Al sumar estas matrices, obtenemos la siguiente forma matricial para la exponencial de la caja de Jordan:
\[
    e^{t J_{\lambda,n}} = e^{\lambda t} 
    \begin{pmatrix}
        1 & 0 & \dots & 0 & 0 \\
        t & 1 & \ddots & \vdots & \vdots \\
        \frac{t^2}{2!} & t & \ddots & 0 & 0 \\
        \vdots & \vdots & \ddots & 1 & 0 \\
        \frac{t^{n-1}}{(n-1)!} & \dots & \frac{t^2}{2!} & t & 1
    \end{pmatrix}
\]

Veamos cómo son las cajas cuando los autovalores son complejos:
Si $\lambda = \alpha + i\beta$ y $\overline{\lambda} = \alpha - i\beta$, en base real aportan un bloque fundamental $J_{\lambda,2}$:
\[
    \begin{pmatrix} \lambda \\ & \overline{\lambda} \end{pmatrix} \xrightarrow{\text{Base Real}} \begin{pmatrix} \alpha & \beta \\ -\beta & \alpha \end{pmatrix}
\]

Queremos poder calcular directamente la exponencial de esta matriz en $\mathbb{R}$ en vez de en $\mathbb{C}$. Para la matriz real de orden 2, tenemos:
\[
    \exp \left[ t \begin{pmatrix} \alpha & \beta \\ -\beta & \alpha \end{pmatrix} \right] = \exp \left[ t\alpha I_2 + t\beta \begin{pmatrix} 0 & 1 \\ -1 & 0 \end{pmatrix} \right] = e^{t\alpha I_2} \cdot \exp \left[ t\beta \begin{pmatrix} 0 & 1 \\ -1 & 0 \end{pmatrix} \right]
\]

Llamemos $K = \begin{pmatrix} 0 & 1 \\ -1 & 0 \end{pmatrix}$. Desarrollando por Taylor la segunda exponencial:
\[
    \exp(t\beta K) = I_2 + t\beta K + \frac{(t\beta)^2}{2!} K^2 + \frac{(t\beta)^3}{3!} K^3 + \frac{(t\beta)^4}{4!} K^4 + \dots
\]

Calculamos las potencias de la matriz $K$:
\begin{align*}
    K^2 &= \begin{pmatrix} 0 & 1 \\ -1 & 0 \end{pmatrix} \begin{pmatrix} 0 & 1 \\ -1 & 0 \end{pmatrix} = \begin{pmatrix} -1 & 0 \\ 0 & -1 \end{pmatrix} = -I_2 \\
    K^3 &= K^2 K = -I_2 K = -K = \begin{pmatrix} 0 & -1 \\ 1 & 0 \end{pmatrix} \\
    K^4 &= K^2 K^2 = (-I_2)(-I_2) = I_2
\end{align*}

Agrupando términos, reconocemos las series de Taylor del seno y el coseno:
\begin{align*}
    \exp(t\beta K) &= I_2 + t\beta K - \frac{(t\beta)^2}{2!} I_2 - \frac{(t\beta)^3}{3!} K + \frac{(t\beta)^4}{4!} I_2 + \dots \\
    &= \left( 1 - \frac{(t\beta)^2}{2!} + \frac{(t\beta)^4}{4!} - \dots \right) I_2 + \left( t\beta - \frac{(t\beta)^3}{3!} + \dots \right) K \\
    &= \cos(\beta t) I_2 + \sin(\beta t) K \\
    &= \begin{pmatrix} \cos(\beta t) & \sin(\beta t) \\ -\sin(\beta t) & \cos(\beta t) \end{pmatrix}
\end{align*}

Por lo tanto, multiplicando por la primera exponencial $e^{t\alpha I_2} = e^{\alpha t} I_2$, llegamos al resultado final:
\[
    \exp \left[ t \begin{pmatrix} \alpha & \beta \\ -\beta & \alpha \end{pmatrix} \right] = e^{\alpha t} \begin{pmatrix} \cos(\beta t) & \sin(\beta t) \\ -\sin(\beta t) & \cos(\beta t) \end{pmatrix}
\]

\textbf{Nota:} Esto es análogo a la fórmula de Euler: $e^{(\alpha + i\beta)t} = e^{\alpha t} (\cos(\beta t) + i \sin(\beta t))$.


\ejemplo{
    Dada la matriz 
    \[
        J = \begin{pmatrix}
            \alpha & 0 & 0 & 0 & 0 & 0 & 0 \\
            0 & \alpha & 0 & 0 & 0 & 0 & 0 \\
            0 & 0 & \beta & 0 & 0 & 0 & 0 \\
            0 & 0 & 1 & \beta & 0 & 0 & 0 \\
            0 & 0 & 0 & 1 & \beta & 0 & 0 \\
            0 & 0 & 0 & 0 & 0 & \gamma & 0 \\
            0 & 0 & 0 & 0 & 0 & 1 & \gamma
        \end{pmatrix}
    \]
    entonces 
    \[
        e^{tJ} = \begin{pmatrix}
            e^{\alpha t} & 0 & 0 & 0 & 0 & 0 & 0 \\
            0 & e^{\alpha t} & 0 & 0 & 0 & 0 & 0 \\
            0 & 0 & e^{\beta t} & 0 & 0 & 0 & 0 \\
            0 & 0 & t e^{\beta t} & e^{\beta t} & 0 & 0 & 0 \\
            0 & 0 & \frac{t^2}{2!} e^{\beta t} & t e^{\beta t} & e^{\beta t} & 0 & 0 \\
            0 & 0 & 0 & 0 & 0 & e^{\gamma t} & 0 \\
            0 & 0 & 0 & 0 & 0 & t e^{\gamma t} & e^{\gamma t}
        \end{pmatrix}
    \]
}

Sea $A \in \mathcal{M}_N(\mathbb{R}), \ a_{ij} \in \mathbb{R}$, tal que $\lambda = \alpha + i\beta$ es un autovalor complejo de $A$ con $\beta > 0$ y $\overline{\lambda} \in \sigma(A)$ su conjugado.

\[
        \begin{pmatrix}
            \alpha & -\beta \\
            \beta & \alpha
        \end{pmatrix} = J_{\mathbb{R}}(\lambda, 1) \xrightarrow{\text{Exponencial}} e^{t J_{\mathbb{R}}(\lambda, 1)} = e^{\alpha t} \begin{pmatrix} \cos(\beta t) & \sin(\beta t) \\ -\sin(\beta t) & \cos(\beta t) \end{pmatrix}
\]

\[
        \begin{pmatrix}
            \lambda & 0 & 0 & 0 \\
            1 & \lambda & 0 & 0 \\
            0 & 0 & \overline{\lambda} & 0 \\
            0 & 0 & 1 & \overline{\lambda}
        \end{pmatrix} \xrightarrow{\text{Base Real}} \begin{pmatrix}
            \alpha & -\beta & 0 & 0 \\
            \beta & \alpha & 0 & 0 \\
            1 & 0 & \alpha & -\beta \\
            0 & 1 & \beta & \alpha
        \end{pmatrix} = J_{\mathbb{R}}(\lambda, 4)
\]

\[
        e^{t \begin{pmatrix}
            J_{\mathbb{R}}(\lambda, 2) & 0_2 \\
            I_2 & J_{\mathbb{R}}(\lambda, 2)
            \end{pmatrix}} = e^{t \begin{pmatrix}
            J_{\mathbb{R}}(\lambda, 2) & 0_2 \\
            0_2 & J_{\mathbb{R}}(\lambda, 2)
            \end{pmatrix} + t \begin{pmatrix}
            0_2 & 0_2 \\
            I_2 & 0_2
            \end{pmatrix}}
\]

\[
        = e^{t \begin{pmatrix}
            J_{\mathbb{R}}(\lambda, 2) & 0_2 \\
            0_2 & J_{\mathbb{R}}(\lambda, 2)
            \end{pmatrix}} \cdot e^{t \begin{pmatrix}
            0_2 & 0_2 \\
            I_2 & 0_2
            \end{pmatrix}} = \begin{pmatrix}
                e^{t J_{\mathbb{R}}(\lambda, 2)} & 0_2 \\
                0_2 & e^{t J_{\mathbb{R}}(\lambda, 2)}
            \end{pmatrix} \cdot \begin{pmatrix}
                I_2 & 0_2 \\
                t I_2 & I_2
            \end{pmatrix} = \begin{pmatrix}
                e^{t J_{\mathbb{R}}(\lambda, 2)} &
                0_2 \\
                t e^{t J_{\mathbb{R}}(\lambda, 2)} & e^{t J_{\mathbb{R}}(\lambda, 2)}
            \end{pmatrix}
\]

\[
        = \begin{pmatrix}
            e^{\alpha t} \cos(\beta t) & e^{\alpha t} \sin(\beta t) & 0 & 0 \\
            -e^{\alpha t} \sin(\beta t) & e^{\alpha t} \cos(\beta t) & 0 & 0 \\
            t e^{\alpha t} \cos(\beta t) & t e^{\alpha t} \sin(\beta t) & e^{\alpha t} \cos(\beta t) & e^{\alpha t} \sin(\beta t) \\
            -t e^{\alpha t} \sin(\beta t) & t e^{\alpha t} \cos(\beta t) & -e^{\alpha t} \sin(\beta t) & e^{\alpha t} \cos(\beta t)
        \end{pmatrix}
\]


\ejemplo{
    Si tenemos dos cajas de Jordan 
    \[
        J_{\lambda,3} = \begin{pmatrix}
            \lambda & 0 & 0 \\
            1 & \lambda & 0 \\
            0 & 1 & \lambda
        \end{pmatrix}
    \]
    con base $\mathcal{B} = \{v_1, v_2, v_3\}$ y
    \[
        J_{\overline{\lambda},3} = \begin{pmatrix}
            \overline{\lambda} & 0 & 0 \\
            1 & \overline{\lambda} & 0 \\
            0 & 1 & \overline{\lambda}
        \end{pmatrix}
    \]
    con base $\overline{\mathcal{B}} = \{\overline{v_1}, \overline{v_2}, \overline{v_3}\}$ 
    entonces la matriz de $A$ restringida al espacio generado por estas cadenas de Jordan con base real es:
        
    \[
        \begin{pmatrix}
            J_{\mathbb{R}}(\lambda, 2) & 0_2 & 0_2 \\
            I_2 & J_{\mathbb{R}}(\lambda, 2) & 0_2 \\
            0_2 & I_2 & J_{\mathbb{R}}(\lambda, 2) 
        \end{pmatrix}  
    \]
    En general, si las cajas son de orden $n$, la matriz de $A$ restringida al espacio generado por estas cadenas de Jordan con base real es:
    \[
        J_{\mathbb{R}}(\lambda, 2n) = \begin{pmatrix}
            J_{\mathbb{R}}(\lambda, 2) & 0_2 & \dots & 0_2 \\
            I_2 & J_{\mathbb{R}}(\lambda, 2) & \dots & 0_2 \\
            \vdots & \ddots & \ddots & \vdots \\
            0_2 & \dots & I_2 & J_{\mathbb{R}}(\lambda, 2)
        \end{pmatrix}
    \]
    Calculamos la exponencial de esta matriz para el caso general. Sea $S$ la matriz diagonal por bloques y $N$ la matriz nilpotente inferior por bloques:
    \[
        e^{t J_{\mathbb{R}}(\lambda, 2n)} = e^{t S + t N} = e^{t S} \cdot e^{t N}
    \]
    Donde:
    \[
        e^{t S} = \begin{pmatrix}
            e^{t J_{\mathbb{R}}(\lambda, 2)} & 0_2 & \dots & 0_2 \\
            0_2 & e^{t J_{\mathbb{R}}(\lambda, 2)} & \dots & 0_2 \\
            \vdots & \vdots & \ddots & \vdots \\
            0_2 & 0_2 & \dots & e^{t J_{\mathbb{R}}(\lambda, 2)}
        \end{pmatrix}, \quad 
        e^{t N} = \sum_{k=0}^{n-1} \frac{t^k N^k}{k!} = \begin{pmatrix}
            I_2 & 0_2 & \dots & 0_2 \\
            t I_2 & I_2 & \dots & 0_2 \\
            \vdots & \ddots & \ddots & \vdots \\
            \frac{t^{n-1}}{(n-1)!} I_2 & \dots & t I_2 & I_2
        \end{pmatrix}
    \]
    Realizando el producto, obtenemos la matriz final:
    \[
        e^{t J_{\mathbb{R}}(\lambda, 2n)} = \begin{pmatrix}
            e^{t J_{\mathbb{R}}(\lambda, 2)} & 0_2 & \dots & 0_2 \\
            t e^{t J_{\mathbb{R}}(\lambda, 2)} & e^{t J_{\mathbb{R}}(\lambda, 2)} & \dots & 0_2 \\
            \vdots & \ddots & \ddots & \vdots \\
            \frac{t^{n-1}}{(n-1)!} e^{t J_{\mathbb{R}}(\lambda, 2)} & \dots & t e^{t J_{\mathbb{R}}(\lambda, 2)} & e^{t J_{\mathbb{R}}(\lambda, 2)}
        \end{pmatrix}
    \]
    donde, si $\lambda = a + ib$, entonces $e^{t J_{\mathbb{R}}(\lambda, 2)} = e^{at} \begin{pmatrix} \cos(bt) & -\sin(bt) \\ \sin(bt) & \cos(bt) \end{pmatrix}$.
}

Sabemos que por el Teorema de Descomposición Primaria, el espacio se descompone en suma directa de los núcleos:
\[
    \mathbb{C}^N = \ker(A - \lambda_1 I)^{\nu(\lambda_1)} \oplus \ker(A - \lambda_2 I)^{\nu(\lambda_2)} \oplus \dots \oplus \ker(A - \lambda_q I)^{\nu(\lambda_q)}
\]
Sea $m_i = m_{alg}(\lambda_i)$ la multiplicidad algebraica de $\lambda_i$ y sean
\[
    \mathcal{B}_1 = \{v_{1,1}, \dots, v_{1,m_1}\} \quad \text{base de} \quad \ker(A - \lambda_1 I)^{\nu(\lambda_1)}
\]
\[
    \mathcal{B}_2 = \{v_{2,1}, \dots, v_{2,m_2}\} \quad \text{base de} \quad \ker(A - \lambda_2 I)^{\nu(\lambda_2)}
\]
\[
    \vdots
\]
\[
    \mathcal{B}_q = \{v_{q,1}, \dots, v_{q,m_q}\} \quad \text{base de} \quad \ker(A - \lambda_q I)^{\nu(\lambda_q)}
\]
entonces $\mathcal{B} = \mathcal{B}_1 \cup \mathcal{B}_2 \cup \dots \cup \mathcal{B}_q$ es una base de $\mathbb{C}^N$.

Consideremos el problema de valor inicial asociado al sistema de ecuaciones diferenciales:
\[
    \begin{cases}
        u' = A u \qquad u(t) = e^{tA} u_0 \\
        u(0) = u_0
    \end{cases}
\]
Descomponemos la condición inicial $u_0$ en la base de la suma directa:
\[
    u_0 = u_{0,1} + u_{0,2} + \dots + u_{0,q}, \quad u_{0,j} \in \ker(A - \lambda_j I)^{\nu(\lambda_j)}, \quad j = 1, \dots, q
\]
Expresando cada término de la suma en función de la base de cada núcleo, tenemos:
\[
    u_0 = \sum_{j=1}^q \sum_{k=1}^{m_j} c_{j,k} v_{j,k}, \quad c_{j,k} \in \mathbb{C}
\]
Para calcular la solución $u(t) = e^{tA} u_0$, aplicamos la linealidad y operamos sobre cada componente $u_{0,j}$:
\[
    u(t) = e^{tA} \left( \sum_{j=1}^q u_{0,j} \right) = \sum_{j=1}^q e^{tA} u_{0,j}
\]
Para cada término de la suma, reescribimos el exponente como $tA = t(\lambda_j I + A - \lambda_j I)$. Como $\lambda_j I$ conmuta con cualquier matriz, podemos separar la exponencial:
\[
    e^{tA} u_{0,j} = e^{\lambda_j t I} e^{t(A - \lambda_j I)} u_{0,j} = e^{\lambda_j t} \left( \sum_{k=0}^{\infty} \frac{t^k}{k!} (A - \lambda_j I)^k \right) u_{0,j}
\] 
Dado que $u_{0,j} \in \ker(A - \lambda_j I)^{\nu(\lambda_j)}$, se cumple que $(A - \lambda_j I)^k u_{0,j} = 0$ para todo $k \geq \nu(\lambda_j)$. Por lo tanto, la serie infinita se trunca y obtenemos la solución exacta en forma finita:
\[
    u(t) = \sum_{j=1}^q e^{\lambda_j t} \left[ \sum_{k=0}^{\nu(\lambda_j)-1} \frac{t^k}{k!} (A - \lambda_j I)^k u_{0,j} \right]
\]
Consideramos ahora el problema anterior pero con la condición inicial $u(0) = v_{j,k}$, es decir, con cada vector de la base de cada núcleo:
\[
    u(t,0,v_{j,k}) = e^{tA} v_{j,k}
\]
entonces $\{e^{tA} v_{j,k} : j = 1, \dots, q, k = 1, \dots, m_j\}$ es una base de soluciones del sistema homogéneo $u' = A u$.
\[
    \Phi(t) = (e^{tA} v_{j,k})_{j=1,\dots,q; k=1,\dots,m_j} = e^{tA} C \quad \text{donde} \quad |C| \neq 0
\]
